% Preamble of report/report.tex.  The mathematical macros are those of theory/preamble.tex,
% so that the proof appendices (report/appendix/A-G) compile unchanged.
\usepackage[utf8]{inputenc}
\usepackage[T1]{fontenc}
\usepackage{lmodern}% required: scalable Latin Modern Type 1 fonts (microtype's font expansion needs scalable fonts; README)
\usepackage[a4paper,margin=2.5cm]{geometry}
\usepackage{amsmath,amssymb,amsthm,mathtools}
\usepackage{enumitem}
\usepackage{booktabs}
\usepackage{array}
\usepackage{tabularx}
\usepackage{longtable}
\usepackage{multirow}
\usepackage{graphicx}
\usepackage{xcolor}
\usepackage{caption}
\usepackage{subcaption}
\usepackage{placeins}
\usepackage[most]{tcolorbox}
\usepackage[numbers,sort&compress]{natbib}
\usepackage{microtype}
\usepackage[hidelinks]{hyperref}
\usepackage[capitalise,nameinlink]{cleveref}

\graphicspath{{figures/}}
\setlength{\emergencystretch}{2em}
\numberwithin{equation}{section}
\captionsetup{font=small,labelfont=bf}

\theoremstyle{plain}
\newtheorem{theorem}{Theorem}[section]
\newtheorem{proposition}[theorem]{Proposition}
\newtheorem{lemma}[theorem]{Lemma}
\newtheorem{corollary}[theorem]{Corollary}
\theoremstyle{definition}
\newtheorem{definition}[theorem]{Definition}
\newtheorem{assumption}[theorem]{Assumption}
\newtheorem{example}[theorem]{Example}
\theoremstyle{remark}
\newtheorem{remark}[theorem]{Remark}

% Claim identifiers (see the introduction of the appendices) and their status labels.
%   S = source result reproved under its own assumptions; D = our derived adaptation;
%   I = identity or definition-level fact;   C = cited for comparison only.
\newcommand{\claim}[2]{\textup{\textsf{[#1\,|\,#2]}}}
% Reference to an item of the source audit (appendix "Source discrepancies and design decisions").
\newcommand{\audit}[1]{\textup{\textsf{[\hyperref[app:audit]{SA\,#1}]}}}
% verified number: \vn{<claim id>}{<literal>} prints the literal; report/verify/claims.py registers each
% claim id with the same literal and the value recomputed from the saved files, and verify_results.py
% checks every occurrence (see appendix J)
\newcommand{\vn}[2]{#2}

\newcommand{\R}{\mathbb{R}}
\newcommand{\N}{\mathbb{N}}
\newcommand{\E}{\mathbb{E}}
\newcommand{\Prob}{\mathbb{P}}
\newcommand{\ind}{\mathbf{1}}
\newcommand{\law}{\mathrm{Law}}
\newcommand{\KL}{\mathrm{KL}}
\newcommand{\JS}{\mathrm{JS}}
\newcommand{\TV}{\mathrm{TV}}
\newcommand{\Ent}{\mathcal{H}}
\newcommand{\cF}{\mathcal{F}}
\newcommand{\cS}{\mathcal{S}}
\newcommand{\cA}{\mathcal{A}}
\newcommand{\cT}{\mathcal{T}}
\newcommand{\dd}{\,\mathrm{d}}
\newcommand{\pre}{\mathrm{pre}}
\newcommand{\SA}{\mathrm{SA}}
\newcommand{\tail}{\mathrm{tail}}
\newcommand{\Cov}{\mathrm{Cov}}
\newcommand{\supp}{\mathrm{supp}}
\DeclareMathOperator*{\argmax}{arg\,max}
\DeclareMathOperator*{\argmin}{arg\,min}

% Symbols that the appendices define with \providecommand, also used in the main text.
\newcommand{\Leb}{\lambda}
\newcommand{\cW}{\mathcal W}
\newcommand{\cB}{\mathcal B}
\newcommand{\Bb}{\mathcal B_b}
\newcommand{\gr}{\mathrm{gr}}

% Method names.
\newcommand{\fPG}{\mbox{$f$-PG}}
\newcommand{\fIRL}{\mbox{$f$-IRL}}
\newcommand{\DMFB}{DM$_{\mathrm{FB}}$}
\newcommand{\ERFB}{ER$_{\mathrm{FB}}$}

% Algorithm boxes (no algorithm package is assumed).
\newtcolorbox{algobox}[1]{breakable,enhanced,colback=white,colframe=black!55,boxrule=0.5pt,arc=1pt,
  left=6pt,right=6pt,top=4pt,bottom=4pt,fonttitle=\bfseries\small,coltitle=black,colbacktitle=black!7,
  title={#1}}
% Answer boxes for the research questions.
\newtcolorbox{answerbox}[1]{breakable,enhanced,colback=black!2,colframe=black!40,boxrule=0.4pt,arc=1pt,
  left=6pt,right=6pt,top=4pt,bottom=4pt,fonttitle=\bfseries\small,coltitle=black,colbacktitle=black!5,
  title={#1}}
